% =====================================================================
%  loop_111_to_100_conversion.tex
%  Drop-in section for the RadCluster formulation paper.
%  Requires: amsmath, amssymb, booktabs, natbib[round,authoryear], hyperref
%  Bibliography: \bibliography{loop_111_to_100_conversion}
% =====================================================================

\section{Burgers-Vector Character and \texorpdfstring{$\tfrac12\langle111\rangle\!\to\!\langle100\rangle$}{1/2<111> to <100>} Loop Conversion}
\label{sec:loop-conversion}

\subsection{Two loop characters and the conversion puzzle}

Irradiation of body-centred-cubic (bcc) iron and ferritic--martensitic steels
produces prismatic interstitial dislocation loops of two Burgers-vector
characters: glissile $\tfrac12 a\langle111\rangle$ loops, which migrate
one-dimensionally along $\langle111\rangle$ and remain mobile to large sizes
($n\gtrsim100$ self-interstitial atoms, SIAs), and sessile $a\langle100\rangle$
loops, which are essentially immobile and adopt a square $\{100\}$ habit
\citep{Masters1965,EyreBartlett1965}.  Isotropic elasticity ($E_l\propto Gb^2$)
favours the smaller Burgers vector $\tfrac12\langle111\rangle$, yet
$a\langle100\rangle$ loops are routinely observed in $\alpha$-Fe and ferritic
alloys, becoming the dominant population at elevated temperature.  Molecular
dynamics shows that displacement cascades inject almost exclusively
$\tfrac12\langle111\rangle$ clusters \citep{Marian2002}; the $\langle100\rangle$
population must therefore be \emph{generated} during microstructural evolution.

We model this generation with two additive channels: a kinetic
collision/junction route \citep{Marian2002} and a thermodynamic
single-loop transformation route \citep{Dudarev2008}.  The former carries the
dose- and size-dependence of $\langle100\rangle$ formation; the latter carries
the temperature trend, in which the $\tfrac12\langle111\rangle$-dominated
microstructure below $\sim350\,^\circ$C gives way to a
$\langle100\rangle$-dominated one above $\sim550\,^\circ$C
\citep{Arakawa2006,Dudarev2008}.

\subsection{Loop free energies}
\label{sec:loop-energies}

The free energy of a prismatic loop of $n$ SIAs and character $X\in\{111,100\}$
is written in the anisotropic-elasticity form of \citet{Dudarev2008}
(equivalently \citealp{Marian2002}, after \citealp{HirthLothe}):
\begin{equation}
  E_l^{X}(n,T) = P_X(n)\!\left[\,
    \hat F_X(T)\,\ln\!\frac{4R^\ast_X(n)}{e\,\delta}
    + F_\delta^{X}(T) + F_c^{X}\,\right],
  \label{eq:loop-energy}
\end{equation}
where $P_X(n)$ is the loop perimeter, $R^\ast_X(n)$ the radius of the equivalent
circular loop, $\delta\simeq0.4$~nm the core cutoff, $\hat F_X$ the
prelogarithmic energy factor (per unit length), $F_\delta^{X}$ the
core-traction term, and $F_c^{X}$ the (zero-temperature) nonlinear core term.
The loop geometry follows from the platelet area $A=n\Omega/b_X$ with $\Omega$
the atomic volume:
\begin{equation}
  R^\ast_X(n)=\sqrt{A/\pi},\quad
  P_{111}(n)=6\sqrt{\tfrac{2A}{3\sqrt3}}\ \text{(hexagon)},\quad
  P_{100}(n)=4\sqrt{A}\ \text{(square)},
  \label{eq:loop-geometry}
\end{equation}
with $b_{111}=\tfrac{\sqrt3}{2}a$ and $b_{100}=a$.  The pure-edge
$\langle100\rangle[100]$ prelogarithmic factor is analytic,
\begin{equation}
  \hat F_{001}([100]) = \frac{a^2}{4\pi}(c_{11}{+}c_{12})
  \!\left[\frac{c_{44}(c_{11}{-}c_{12})}{c_{11}(c_{11}{+}c_{12}{+}2c_{44})}\right]^{1/2}
  \;\xrightarrow{\;c'\to0\;}\;0,
  \label{eq:F100}
\end{equation}
and \emph{vanishes} as the shear constant $c'=\tfrac12(c_{11}-c_{12})$ softens,
whereas all $\tfrac12\langle111\rangle$ factors remain finite.  The softening is
driven by spin fluctuations approaching the $\alpha$--$\gamma$ transition
\citep{Dudarev2008},
\begin{equation}
  c'(T)\simeq 56.8\,(1-T/T_c)^{1/2}\ \text{GPa},\qquad T_c=1185\ \text{K}.
  \label{eq:cprime}
\end{equation}
This is the microscopic origin of the temperature trend: the
$\langle100\rangle$ elastic self-energy collapses at high $T$ while the
$\tfrac12\langle111\rangle$ energy does not.  Representative zero-temperature
core constants from \citet{Dudarev2008} are listed in
Table~\ref{tab:core-constants}.

\subsection{Channel A: thermodynamic single-loop transformation}
\label{sec:channel-A}

A single $\tfrac12\langle111\rangle$ loop of size $n$ reorients to
$\langle100\rangle$ when the latter becomes the lower-energy configuration.  The
driving force is the per-loop free-energy difference
\begin{equation}
  \Delta F(n,T) = E_l^{111}(n,T) - E_l^{100}(n,T) = P(n)\,\Delta f(T),
  \label{eq:driving-force}
\end{equation}
which we parametrise through a per-unit-length difference $\Delta f(T)$ tied to
the same softening law~\eqref{eq:cprime},
\begin{equation}
  \boxed{\;\Delta f(T) = f_{111} - f_{100}^{0}\,(1-T/T_c)^{p}\;}
  \label{eq:delta-f}
\end{equation}
with $f_{111},f_{100}^{0}$ the (near-zero-$T$) per-length energies of the two
characters and the exponent $p$ fixed so that $\Delta f$ changes sign near
$350\,^\circ$C and is strongly positive by $550\,^\circ$C, reproducing the three
stability regions of \citet{Dudarev2008}.  Conversion is thermodynamically
admissible for $T>T^\ast$, where $f_{111}=f_{100}^{0}(1-T^\ast/T_c)^{p}$.

The transformation is a thermally activated, one-way process.  Because the
reorientation propagates segment-by-segment along the loop
\citep{Marian2002}, its effective barrier scales with the number of dislocation
segments to be reoriented, i.e.\ with the perimeter:
\begin{equation}
  \boxed{\;\Gamma_{\rm uni}(n,T) = \nu_0\,
    \exp\!\left[-\frac{E_a^{\rm uni}(n)}{k_BT}\right]
    \max\!\left[0,\;1-\exp\!\left(-\frac{\Delta F(n,T)}{k_BT}\right)\right],\qquad
    E_a^{\rm uni}(n) = E_a^{0} + \gamma_a\,\frac{P(n)}{b_{111}}\;}
  \label{eq:gamma-uni}
\end{equation}
where $\nu_0\sim10^{13}\,\mathrm{s}^{-1}$ is an attempt frequency, $E_a^{0}$ an
offset barrier ($\sim$ the rate-limiting $\tfrac12\langle110\rangle\!\to\!\langle100\rangle$
step, $\sim1$~eV, \citealp{Marian2002}), and $\gamma_a$ sets the size
suppression.  The gating factor favours large loops (larger $\Delta F$) while the
size-dependent barrier suppresses them, producing a preferred conversion-size
window: this reproduces the spontaneous transformation of \emph{small} loops on
heating \citep{Arakawa2006} while leaving large mobile $\tfrac12\langle111\rangle$
loops to convert through Channel~B.

\subsection{Channel B: kinetic junction nucleation and absorption growth}
\label{sec:channel-B}

\paragraph{Nucleation.}  Two mobile $\tfrac12\langle111\rangle$ loops of
\emph{comparable} size that collide can react through the junction
\citep{EyreBullough1965,Marian2002}
\begin{equation}
  \tfrac12[111] + \tfrac12[1\bar1\bar1] \rightarrow [100],
  \label{eq:junction}
\end{equation}
condensing into a single $\langle100\rangle$ loop of the combined SIA content,
driven by the energy reduction of merging.  Molecular dynamics shows the
junction forms only when the two loops are of similar size; otherwise the
smaller simply rotates into the orientation of the larger (ordinary
coalescence).  We therefore split the $\tfrac12\langle111\rangle$ collision
kernel $\mathcal K^{ii}_{n,n'}$ (Eq.~\ref{eq:Kii} below) with a
size-comparability branching fraction
\begin{equation}
  \boxed{\;K^{\rm junc}_{n,n'} = \varphi_{\rm junc}(n,n')\,\mathcal K^{ii}_{n,n'},
  \qquad
  \varphi_{\rm junc}(n,n') = \varphi_{\max}
  \exp\!\left[-\frac{\big(\ln(n/n')\big)^2}{2\sigma_s^2}\right]
  \Theta\!\big(\min(n,n')\ge n_{j,\min}\big)\;}
  \label{eq:phi-junc}
\end{equation}
the product ($\langle100\rangle$, size $n{+}n'$) entering the
$\langle100\rangle$ population and the complement $(1-\varphi_{\rm junc})$
remaining ordinary $\tfrac12\langle111\rangle$ coalescence.  Here
$\varphi_{\max}\!\le\!1$ is the peak yield at $n=n'$, $\sigma_s$ the log-size
tolerance, and $n_{j,\min}$ the minimum junction size.  The same-polarity
collision kernel is the standard size-resolved form
\begin{equation}
  \mathcal K^{ii}_{n,n'} = \frac{8\pi}{\Omega^{2/3}}
  \big(\xi_n+\xi_{n'}\big)\big(D^{111}_n+D^{111}_{n'}\big),\qquad
  \xi_n=\Big(\tfrac{3n}{8\pi}\Big)^{1/3}.
  \label{eq:Kii}
\end{equation}

\paragraph{Growth by absorption.}  Once formed, $\langle100\rangle\{100\}$ loops
are effectively stationary (glide barrier $>2.5$~eV) and act as biased sinks for
mobile cascade-produced $\tfrac12\langle111\rangle$ clusters \citep{Marian2002},
\begin{equation}
  [100]_m + \tfrac12\langle111\rangle_n \rightarrow [100]_{m+n},
  \label{eq:absorption}
\end{equation}
the absorbed cluster rotating into the $\langle100\rangle$ orientation.  This is
the dominant route by which $\langle100\rangle$ loops reach TEM-visible sizes.
The capture kernel is the cross-character collision rate, driven by the mobile
$\tfrac12\langle111\rangle$ partner ($D^{100}_m\approx0$):
\begin{equation}
  \boxed{\;K^{\rm abs}_{m,n} = \frac{8\pi}{\Omega^{2/3}}
  \big(\xi_m+\xi_n\big)\big(D^{100}_m+D^{111}_n\big)
  \approx \frac{8\pi}{\Omega^{2/3}}\big(\xi_m+\xi_n\big)D^{111}_n\;}
  \label{eq:Kabs}
\end{equation}

\subsection{Incorporation into the cluster-dynamics equations}
\label{sec:cd-incorporation}

The single SIA cluster population $c_n$ is split into two coupled populations:
$c_n^{(111)}$ (glissile, mobile to $n_{\max}^i$, fed by all cascade SIA
production, undergoing self-coalescence) and $c_m^{(100)}$ (sessile,
$m\ge n_{\rm conv}$, no cascade source, no self-coalescence).  The
$\langle100\rangle$ onset size $n_{\rm conv}$ is fixed self-consistently as the
smallest $n$ with $\Delta F(n,T)>0$.  The added terms in the master equation are

\begin{align}
  \frac{dc_n^{(111)}}{dt}\Big|_{\rm conv}
    &= -\,\Gamma_{\rm uni}(n,T)\,c_n^{(111)}
       \;-\; c_n^{(111)}\!\sum_{n'} K^{\rm junc}_{n,n'}\,c_{n'}^{(111)}
       \;-\; c_n^{(111)}\!\sum_{m} K^{\rm abs}_{m,n}\,c_m^{(100)},
       \label{eq:me-111}\\[4pt]
  \frac{dc_m^{(100)}}{dt}\Big|_{\rm conv}
    &= +\,\Gamma_{\rm uni}(m,T)\,c_m^{(111)}
       \;+\; \tfrac12\!\!\sum_{n+n'=m}\!\! K^{\rm junc}_{n,n'}\,c_n^{(111)}c_{n'}^{(111)}
       \;+\!\!\sum_{n}\! K^{\rm abs}_{m-n,n}\,c_{m-n}^{(100)}c_n^{(111)}
       \;-\!\!\sum_{n}\! K^{\rm abs}_{m,n}\,c_m^{(100)}c_n^{(111)}.
       \label{eq:me-100}
\end{align}

In addition, $c_m^{(100)}$ obeys the usual point-defect absorption (SIA and
vacancy monomer capture), thermal emission with a $\langle100\rangle$-specific
binding energy $E_b^{100}(m)=A_{100}m^{-B_{100}}$, fixed-sink loss, and
vacancy--SIA recombination/annihilation, all in their sessile form.  Every term
in \eqref{eq:me-111}--\eqref{eq:me-100} conserves the SIA content
$\sum_n n\,c_n^{(111)}+\sum_m m\,c_m^{(100)}$ (the junction and absorption
products carry the summed sizes; the unary transfer is size-preserving), so the
Frenkel-pair conservation diagnostic is unaffected.

\subsection{Input data and parameters}
\label{sec:loop-conversion-data}

\begin{table}[ht]
\centering
\caption{Zero-temperature core constants for the loop free
energy~\eqref{eq:loop-energy}, after \citet{Dudarev2008} (eV/\AA).}
\label{tab:core-constants}
\begin{tabular}{lcccc}
\toprule
term & $111[11\bar2]$ & $111[\bar110]$ & $001[100]$ & $001[110]$ \\
\midrule
$F_c$ (nonlinear core)   & 0.46 & 0.47 & 0.33 & --- \\
$F_\delta$ (core-traction) & 0.345 & 0.349 & 0.387 & 0.390 \\
\bottomrule
\end{tabular}
\end{table}

\begin{table}[ht]
\centering
\caption{Conversion-model parameters.  Values marked ``cal.'' are calibrated
against the measured $\tfrac12\langle111\rangle$ fraction $f_{111}(T,\text{dose})$.}
\label{tab:conversion-params}
\begin{tabular}{llll}
\toprule
symbol & meaning & first-cut value & source \\
\midrule
$T_c$            & $\alpha$--$\gamma$ softening temperature & 1185 K & \citet{Dudarev2008} \\
$c'(T)$          & shear softening law & $56.8(1{-}T/T_c)^{1/2}$ GPa & \citet{Dudarev2008} \\
$\delta$         & dislocation core cutoff & $0.4$ nm & \citet{Dudarev2008} \\
$f_{111},f_{100}^{0}$ & per-length loop energies & $0.3$--$0.5$ eV/\AA & Tab.~\ref{tab:core-constants} \\
$p$              & softening exponent in $\Delta f(T)$ & cal.\ ($\Delta f{=}0$ near $350\,^\circ$C) & \citet{Dudarev2008} \\
$\nu_0$          & attempt frequency & $10^{13}\,\mathrm{s}^{-1}$ & Debye \\
$E_a^{0}$        & unary barrier offset & $0.5$--$1.0$ eV & \citet{Marian2002} \\
$\gamma_a$       & unary barrier size slope & cal. & \citet{Arakawa2006} \\
$\varphi_{\max}$ & peak junction yield ($n{=}n'$) & $0.5$--$1.0$ & \citet{Marian2002} \\
$\sigma_s$       & log-size tolerance & $0.3$--$0.5$ & \citet{Marian2002} \\
$n_{j,\min}$     & minimum junction size & $\mathcal O(10)$ SIAs & \citet{Marian2002} \\
$A_{100},B_{100}$& $\langle100\rangle$ binding-energy fit & from Eq.~\eqref{eq:loop-energy} & this work \\
$n_{\rm conv}$   & $\langle100\rangle$ onset size & \emph{derived} ($\Delta F{>}0$) & Eq.~\eqref{eq:driving-force} \\
\bottomrule
\end{tabular}
\end{table}

% \bibliographystyle{plainnat}
% \bibliography{loop_111_to_100_conversion}
